LLM-based code-review agents are increasingly entrusted with gating code
changes, yet current benchmarks largely measure whether an agent produces the
correct accept/reject label without separating two distinct questions: whether
the \emph{evidence} presented is intact and trustworthy, and whether it is
\emph{sufficient} to justify a decision at all. We introduce
\textbf{GroundedReview}, a benchmark that (i) casts code review as a three-way
decision---\ACCEPT, \REJECT, or \NEED---and (ii) evaluates the same real
issue--patch pair under several evidence conditions that independently control
which test outcomes are visible, whether the agent may actively fetch
artifacts, and whether provenance and freshness are verifiable. We distinguish
\emph{Oracle Gold} (the decision warranted by a fresh, scoped pytest run) from
\emph{Condition Gold} (the decision warranted by the evidence actually
granted, which is legitimately \NEED\ on restricted surfaces). On ten SymPy
issue--patch pairs served to two models (\texttt{deepseek-flash} and
\texttt{gpt-4.1-mini}) across five conditions, a corrective replication shows
that repairing protocol-construct defects in our own instrument moved one
model's sufficient-evidence abstention from $9/10$ to $1/10$---an
associational, not causal, contrast---whereas the other model, under a
restricted ``strict'' surface whose correct answer is \NEED, returned
unsupported \ACCEPT\ on $7/10$ issues, upgrading non-decisive evidence into
acceptance. A traceable case shows an agent accepting a patch after an
already-visible fail-to-pass test had failed. A disjoint eight-issue,
80-cell capped exploratory run repeats the model asymmetry associationaly.
We then run a \emph{pre-registered controlled intervention} (G4) on 18
frozen issue--patch pairs: a yoked, fixed nine-step evidence transcript in a
$2\times2$ design (144 cells, one call each, US\,\$0.807 of a US\,\$2.00 hard
cap) that holds linkage constant and independently varies decisive-outcome
visibility and freshness/coverage attestation. With decisive results blocked,
adding integrity attestation moves GPT-4.1-mini from $0/18$ to $8/18$
unsupported acceptance, with no reverse moves (paired two-sided exact McNemar
$p=0.0078$; paired difference $-0.444$, issue-clustered 95\,\% CI
$[-0.667,-0.222]$), including acceptance of two of five oracle-\REJECT\
patches; DeepSeek-Flash does not move ($18/18$ \NEED\ in both cells,
$p=1.0$). When decisive outcomes are visible, both models reject all five
failing patches (oracle agreement 16--17/18). The effect is therefore real
under the frozen protocol but model-specific: one model treats intact
integrity attestations as a warrant to accept despite withheld decisive
results, while the other does not. An independent third-party human second
annotation of all 18 blind packets has since been obtained: blind verdicts
agree with the provisional oracle on 17/18 issues, the two focused reviews
close FR1 as an instrument-only fix and FR2 as a single REJECT standard,
and a human arbiter overrides issue~19346 to \REJECT\ (one disagreement on
issue~15875 remains unresolved). Labels are still reported as
\textbf{PROVISIONAL} until the signed attestation is filed; we therefore
present a measurement framework and clearly bounded evidence rather than
population-level claims about model behaviour.
